Source : Eurostat, comptes annuels des administrations publiques (gov\_10a\_main, gov\_10a\_taxag, gov\_10a\_exp), PIB (nama\_10\_gdp) ; Insee, dette trimestrielle de Maastricht par sous-secteur (contributions consolidées, 4e trimestre). Calculs de l'auteur (dette/fongible.py).\par Lecture : la dépense de vieillesse est passée de 11,0 \% à 13,4 \% du PIB. Hors intérêts, la dépense publique a augmenté de 2,3 points ; la santé et la vieillesse en expliquent 4,4, les autres fonctions ont baissé ensemble de 2,1 points.
